\subsection*{12.3 Orthogonality Principle}

When W has infinite dimension, computing the projection of a random variable

Y onto a closed space W using the definition may be difficult. The next theorem

provides a useful characterization of the projection function in terms of perpen-

dicular projection. Essentially, the theorem states that the optimal solution in W

that minimizes the distance to Y is achieved when and only when the error vector

is perpendicular to W Note that we are not proving any existence result in the

theorem below, as the well-definedness of the projection function has already been

established in Theorem 12.4. The objective of the theorem below is to give an

alternate way to compute the projection of a random variable.

\begin{thmbox}{12.5 (Orthogonality Principle)}
\end{thmbox}

Given a random variable Y in L2(P) and a closed subspace W in L2(P) ,t h e

projection of Y onto W is characterized by

X* =Proj W (Y) X * \inW and \langleY- X *,X \rangle= 0 X\in W.

\textit{Proof} ( . ) Suppose. X* is the random variable in W that minimizes Y- X over

all X\in W We need to show that \langleY- X *,X \rangle= 0 for all X\in W We proceed by

contradiction. Suppose there exists a random variable X0 in W such that Y- X * is

not orthogonal to X0, ie.,\langleY- X *,X 0\rangle /=0 We can then deduce that X* is not the

minimizer.

One way to see this by completing the square. Consider

Y- (X * -uX 0)2

as a function of u Since X* minimizes Y- X over all X\in W , this function is

minimized atu= 0 Expanding it as a polynomial in u ,we have

uX0 +Y - X *2 =u 2X02 +2 u\langleX0,Y - X *\rangle+ Y- X *2.

Since it is assumed that \langleX0,Y - X *\rangle /=0 ,t h e. L2 normX02 should be nonzero.

We complete the square and write the above expression as

X02

(

u+ \langleX0,Y - X *\rangle

X02

) 2

- \langleX0,Y - X *\rangle2

X02 + Y- X *2.

We can pick the value of u such that the square in the first term is zero. Then,

- \langleX0,Y - X *\rangle2

X02 + Y- X *2

is strictly less than Y- X *2. This contradicts the minimality of Y- X *2, and

hence, we must have \langleY- X *,X 0\rangle= 0.

(. ) Suppose. X is a random variable such that

\langleY- X,X \rangle= 0 for allX \in W. (12.2)

We want to show that X coincides with ProjW (Y) We havej u s tp r o v e di nt h e

previous paragraph that ProjW (Y)= X * satisfies

\langleY- Proj W (Y), X\rangle= 0 for allX \in W. (12.3)

By subtracting (12.3) from (12.2), we get

\langleProjW (Y)- X,X \rangle= 0 for allX \in W. (12.4)

Moreover,ProjW (Y)- X is a random variable in W because both ProjW (Y) and. X

are in W and W is a subspace. Hence, in (12.4), we can substitute X by. ProjW (Y)-

X to getProj W (Y)- X2 =0 This proves that ProjW (Y)= X almost surely. \hfill $\square$

In view of the Orthogonality Principle, we can interpret the minimum mean-

squared error (MMSE) estimator geometrically as a projection operator.
